\section{Claim Checking and Bounded Repair}
The post-processor splits generated text into factual spans while retaining character offsets and treating citation markers as separate spans. It resolves cited evidence IDs to source pages and records whether a citation came from the model, mapping, or repair. A deterministic token-overlap score compares each claim with its cited source text; numeric mismatch and threshold rules assign operational labels. These labels are \emph{lexical routing decisions}, not validated entailment or contradiction judgments. Pixel-only claims remain explicitly unverified. A negative-control test of the lexical support rule failed, so we do not use its labels as the paper's correctness endpoint.

Repair uses a bounded action set: keep a span, narrow a partial clause, attempt one citation remap, delete a span, or abstain if no claim remains. It does not generate replacement facts. A second pass records the resulting spans and edits, but reusing the same lexical rule is not independent confirmation. The paired study isolates this post-processing intervention by applying both settings to the same raw completion and prompt. It asks whether the rule changes observable overlap, page alignment, and status; whether humans find the repaired answer better is unmeasured.

This separation between an unchanged raw answer and a bounded post-processor determines the comparisons and denominators in the evaluation design.
